Um zu zeigen, dass die inverse Transformation $S^{-1}(\Lambda)$ als $\left(\gamma^0 S(\Lambda) \gamma^0\right)^{\dagger}$ geschrieben werden kann, beginnen wir mit der Definition von $S(\Lambda)$:

\[
S(\Lambda) = e^{-\frac{i}{4} \omega^{\mu \nu} \sigma_{\mu \nu}}
\]

wobei $\sigma_{\mu \nu} = \frac{i}{2} [\gamma_\mu, \gamma_\nu]$ ist. Die Inverse von $S(\Lambda)$ ist gegeben durch:

\[
S^{-1}(\Lambda) = e^{\frac{i}{4} \omega^{\mu \nu} \sigma_{\mu \nu}}
\]

Um die Beziehung $S^{-1}(\Lambda) = \left(\gamma^0 S(\Lambda) \gamma^0\right)^{\dagger}$ zu zeigen, betrachten wir die Eigenschaften der Dirac-Matrizen und die Struktur von $\sigma_{\mu \nu}$. Die Dirac-Matrizen erfüllen die Beziehung:

\[
(\gamma^0 \gamma^\mu \gamma^0) = (\gamma^\mu)^\dagger
\]

Daraus folgt, dass:

\[
(\gamma^0 \sigma_{\mu \nu} \gamma^0) = \sigma_{\mu \nu}^\dagger
\]

Da $S(\Lambda)$ exponentiell in $\sigma_{\mu \nu}$ ist, ergibt sich:

\[
\left(\gamma^0 S(\Lambda) \gamma^0\right) = e^{-\frac{i}{4} \omega^{\mu \nu} (\gamma^0 \sigma_{\mu \nu} \gamma^0)} = e^{-\frac{i}{4} \omega^{\mu \nu} \sigma_{\mu \nu}^\dagger}
\]

Das bedeutet, dass:

\[
\left(\gamma^0 S(\Lambda) \gamma^0\right)^{\dagger} = \left(e^{-\frac{i}{4} \omega^{\mu \nu} \sigma_{\mu \nu}^\dagger}\right)^{\dagger} = e^{\frac{i}{4} \omega^{\mu \nu} \sigma_{\mu \nu}}
\]

Dies zeigt, dass:

\[
S^{-1}(\Lambda) = \left(\gamma^0 S(\Lambda) \gamma^0\right)^{\dagger}
\]

Diese Beziehung folgt aus der speziellen Struktur der Dirac-Matrizen und der Definition der Lorentz-Transformation für Dirac-Spinoren.